\documentclass[10pt]{article}
\input{../../notes-style.tex}

\usepackage{multicol}
\usepackage{array}

\title{Ch.~2 Cheat Sheet\\\large Algorithms, Recursion, Strategies}
\author{}
\date{}

\begin{document}
\maketitle
\vspace{-2em}

\begin{multicols}{2}

\subsection*{Problem vs.\ algorithm}
A \textbf{problem} specifies inputs $\to$ outputs (\emph{what}). An
\textbf{algorithm} is a procedure that solves it (\emph{how}). One problem
usually has many correct algorithms; you pick by correctness, speed,
memory, simplicity.

\subsection*{What makes an algorithm ``good''}
\begin{enumerate}
    \item \textbf{Correct} on all inputs (including edges).
    \item \textbf{Fast enough} for input sizes you'll see.
    \item \textbf{Reasonable memory}.
    \item \textbf{Simple enough} to implement, test, maintain.
\end{enumerate}

\subsection*{Recursion anatomy}
Every recursive function must have:
\begin{itemize}
    \item A \textbf{base case} returning without recursing.
    \item A \textbf{recursive case} that calls itself on a
    \emph{strictly smaller} input.
\end{itemize}
If the recursive call isn't smaller or doesn't converge to the base, you
stack-overflow.

\begin{lstlisting}[language=C++]
int factorial(int n) {
    if (n <= 1) return 1;        // base
    return n * factorial(n - 1); // recurse
}
\end{lstlisting}

\subsection*{Binary search (recursive)}
\begin{lstlisting}[language=C++]
int bsearch(const vector<int>& v,
            int lo, int hi, int t) {
    if (lo > hi) return -1;      // base: miss
    int m = lo + (hi - lo) / 2;
    if (v[m] == t) return m;
    if (v[m] <  t) return bsearch(v, m+1, hi, t);
    return               bsearch(v, lo, m-1, t);
}
\end{lstlisting}
\textbf{Precondition:} \texttt{v} sorted. \textbf{Cost:} $O(\log n)$.

\subsection*{Recursion gotchas}
\begin{itemize}
    \item Missing base $\to$ infinite recursion / stack overflow.
    \item Recursive call not smaller $\to$ same.
    \item \textbf{Overlapping subproblems} (na\"ive Fib) $\to$
    exponential blowup. Memoize or iterate.
    \item Deep stacks cost memory; iterative form may be required for
    huge $n$.
\end{itemize}

\subsection*{Heuristic vs.\ greedy vs.\ DP}
\begin{tabular}{@{}lp{3.5cm}@{}}
\textbf{Heuristic} & Good-enough rule; no optimality guarantee.\\
\textbf{Greedy}    & Locally optimal choice each step; optimal \emph{only}
when greedy-choice property holds.\\
\textbf{DP}        & Optimal by exploring smaller subproblems, storing
results; needs optimal substructure + overlapping subproblems.
\end{tabular}

\columnbreak

\subsection*{Greedy: when it works}
Use greedy iff the problem has the \textbf{greedy-choice property}: the
locally optimal choice is always part of \emph{some} globally optimal
solution. Classic wins:
\begin{itemize}
    \item Fractional knapsack (sort by value/weight, take).
    \item Activity selection (earliest finish time).
    \item Coin change with ``canonical'' coin systems (US coins OK;
    arbitrary systems \emph{not} OK).
\end{itemize}
Classic fails:
\begin{itemize}
    \item 0/1 knapsack (greedy on ratio is wrong).
    \item Coin change with $\{1, 3, 4\}$, target 6: greedy picks $4{+}1{+}1$,
    optimum is $3{+}3$.
\end{itemize}

\subsection*{DP: the template}
\begin{enumerate}
    \item Identify \textbf{state}: what parameters uniquely determine a
    subproblem?
    \item Write the \textbf{recurrence}: optimal for state $s$ in terms
    of smaller states.
    \item Identify \textbf{base cases}.
    \item Choose \textbf{order} to fill the table (bottom-up) or use
    memoization (top-down).
    \item Space-optimize if only the last row/col is needed.
\end{enumerate}
Signals a problem is DP: asks for optimum/count/exists; recurrence
reuses same subproblem many times.

\subsection*{Recognizing the strategy}
\begin{tabular}{@{}lp{4cm}@{}}
``Find optimum'' + choice-at-each-step, greedy works &
\textbf{Greedy}\\
``Find optimum'' + overlapping subproblems   & \textbf{DP}\\
Too expensive to solve exactly, need ``good enough'' & \textbf{Heuristic}\\
Natural smaller-subproblem structure, each solved once & \textbf{Div \& Conq}
\end{tabular}

\subsection*{Data privacy \& ethics -- the two levers}
\begin{itemize}
    \item \textbf{Collection.} Only gather what you need; de-identify
    when possible; plan retention.
    \item \textbf{Release.} Aggregate, suppress small cells, use
    differential privacy or k-anonymity; assume auxiliary data exists.
\end{itemize}
Re-identification via auxiliary data is the most common failure mode
(Netflix, AOL, NYC taxi logs).

\subsection*{Ethics checklist (CS 300 flavor)}
\begin{itemize}
    \item \textbf{Harm.} Who might this code hurt?
    \item \textbf{Consent.} Did users agree to \emph{this} use?
    \item \textbf{Transparency.} Can I explain the decision logic?
    \item \textbf{Fairness.} Does it disparately affect a group?
    \item \textbf{Reproducibility.} Could someone rerun this and audit?
\end{itemize}

\subsection*{Top gotchas}
\begin{itemize}
    \item Claiming greedy is optimal without checking the
    greedy-choice property.
    \item Writing DP without defining the state first.
    \item Na\"ive recursion on a problem with overlapping subproblems
    (Fibonacci trap).
    \item Using a heuristic result as if it were optimal.
\end{itemize}

\end{multicols}
\end{document}
